\section{Why simultaneous local exceptions have product size}\label{sec:joint-rank}

\status{The passage from single-place tests to the joint exceptional-generator bound is proved here. Producing the quantitative single-place tests, and the integral polynomial lattice comparisons used below, remain explicit inputs. This section does not repeat the conversion of elementary-divisor tails to index.}

Work over a DVR $\mathcal O$ with fraction field $K$, maximal ideal $\mathfrak m$, and a nonzero nonunit $p$; in the source $w^2=p$. Fix an integer $D\ge0$. The base and dual coefficient lattices are finite free of constant rank $d_v$. At a target lattice $L$, the relevant dual-slot lattice is
\[
 P_v(L)=N_{\mathrm{tgt},v}(L)^\vee
   =\{x:\ell(x)\in\mathcal O\text{ for every }\ell\in N_{\mathrm{tgt},v}(L)\},
 \qquad \operatorname{rank}P_v(L)=d_v.
\]
The local input supplies finitely many incoming edges $e_i=(L_i^-,L)$ and vectors $F_i$ in the source lattice at $L_i^-$. Put $f_i=p^{A_v}F_i$, where $A_v$ is the source's sum of the first $r$ weights. The integral pairing statements imply $f_i\in P_v(L)$. Besides vanishing to order $w^{c_*}$ at every other target of $L_i^-$, the input supplies $n_v\ge(1-\kappa_v)d_v$, an $\mathcal O$-linear map $x\mapsto\ell_x$ into $N_{\mathrm{tgt},v}(L)$, and an integral matrix $C$ such that
\begin{equation}\label{eq:jr-tests}
 \sum_i C_{ji}\ell_x(f_i)=p^D x_j
 \qquad(x\in\mathcal O^{n_v},\ 1\le j\le n_v).
\end{equation}
Here $\kappa_v\ge0$ and $c_*\ge4(D+1)$. This is the concrete rank content of \texttt{LPprime}, rather than an assumed independence of exceptional subspaces. \textsf{[JR1, JR12]}

\begin{lemma}[Extracting good integral directions]\label{lem:jr-smith}
Under \eqref{eq:jr-tests}, $P_v(L)$ has a basis $b_{v,a}$ with at most $\kappa_vd_v$ bad indices, such that
\[
 p^D b_{v,a}\in\sum_i\mathcal O f_i
 \quad\text{for every good index }a.
\]
\end{lemma}
\begin{proof}
Put $q=p^D$ and apply Smith normal form to $\sum_i\mathcal O f_i\subseteq P_v(L)$, with nonzero elementary divisors $a_i$. The integral functionals $\ell_x$ and \eqref{eq:jr-tests} give an integral matrix identity $C'\operatorname{diag}(a_i)X=qI_{n_v}$. Call $i$ good if $a_i\mid q$. For every other $i$, valuation comparability gives $a_i=q\mu_i$ with $\mu_i\in\mathfrak m$. Moving these terms to the right gives
\[
 \sum_{i\ \mathrm{good}} C'_i a_i X_i=q(I-Y),
 \qquad Y\in M_{n_v}(\mathfrak m).
\]
The right side has rank $n_v$ over $K$, since $I-Y$ has unit determinant. The left side is a sum of at most one rank-one matrix per good index. Hence there are at least $n_v$ good indices. For their ambient Smith basis vectors, $q b_{v,a}$ lies in the generated sublattice. The other at most $d_v-n_v\le\kappa_vd_v$ indices are bad. \textsf{[JR4]}
\end{proof}

To globalize, use the finite set $\Omega$ of orbits of decorated lattice tuples. The acting groups have trivial stabilizers and preserve the coefficient lattices. Consequently a value $\xi$ at a representative $L$ extends uniquely to an equivariant section $\delta_L(\xi)$ supported on that orbit. Let $M_B$ use the dual coefficient lattice at coordinates $v\in B$ and the base coefficient lattice elsewhere. Let $J_B\subseteq M_B$ be the sum, over $v\in B$, of the neighbor-operator images with the source lattice in coordinate $v$. The neighbor operator itself includes the factor $p^{A_v}$; $A_v$ may be negative.

For a good index at $v$, tensor its relation in the lemma with arbitrary integral vectors in the other coordinates. Applying the neighbor operator to the corresponding section at $L_i^-$ produces the desired section at $L$ and terms at the other neighbors. The desired term occurs once, because stabilizers are trivial. All other terms lie in $w^{c_*}M_B\subseteq p^{D+1}M_B$. Equivariance gives the same assertion at every translate. Therefore every elementary tensor with at least one good selected coordinate satisfies
\begin{equation}\label{eq:jr-one-good}
 p^D\delta_L(b_{1,a_1}\otimes\cdots\otimes b_{s,a_s})
       \in J_B+p^{D+1}M_B.
\end{equation}
This cancellation uses the integral neighbor-map inclusion and equivariance as hypotheses; their actual polynomial verification is retained in \textsf{[JR12]}. The cancellation itself is proved by the preceding evaluation argument. \textsf{[JR5]}

\begin{proposition}[Joint exceptional generators]\label{prop:jr-joint}
Suppose these local tests and lattice properties hold for every selected coordinate, with common $D,c_*$. Write $d=\operatorname{rank}M_{\mathrm{na}}$ for the full base function lattice. For every nonempty set $B$ of selected coordinates, there is a finite set $\mathcal C_B\subseteq M_B$ with
\begin{equation}\label{eq:jr-joint}
 |\mathcal C_B|\le d\prod_{v\in B}\kappa_v,
 \qquad
 p^{D+1}M_B\subseteq J_B+\operatorname{span}_{\mathcal O}\mathcal C_B.
\end{equation}
\end{proposition}
\begin{proof}
At each orbit representative, retain precisely the elementary tensors whose index is bad at \emph{every} coordinate of $B$. Their number is at most
\[
 \left(\prod_{v\in B}\kappa_vd_v\right)
 \left(\prod_{v\notin B}d_v\right)
   =\left(\prod_{v\in B}\kappa_v\right)\prod_v d_v.
\]
Every other elementary tensor satisfies \eqref{eq:jr-one-good}. Delta-extend the retained tensors and take their union over $\Omega$ to obtain $\mathcal C_B$. The base coefficient tensors, similarly extended, are linearly independent: evaluate a relation at each representative. Thus $|\Omega|\prod_vd_v\le d$, proving the cardinality bound. \textsf{[JR6, JR7]}

With $T=J_B+\operatorname{span}_{\mathcal O}\mathcal C_B$, expansion in the tensor bases gives $p^DM_B\subseteq T+p^{D+1}M_B$. The module
\[
 Q=(p^DM_B+T)/T
\]
is finitely generated and satisfies $Q=pQ$. Nakayama's lemma gives $Q=0$. Multiplying the resulting inclusion by $p$ gives exactly \eqref{eq:jr-joint}, the exponent advertised by the source. \textsf{[JR7]}
\end{proof}

The product appears because a tensor survives only when \emph{all} selected coordinates are bad. Arbitrary subspace intersections do not satisfy such a product bound. Summing the displayed products over $|B|=k$ gives the elementary symmetric polynomial $e_k(\kappa)$; the subsequent lattice walk supplies the additional multiplicities and the exponent of $p$.

\paragraph{What remains at one place.}
The supplementary report proves how primitive pairings isolated at different incoming edges combine to give \eqref{eq:jr-tests}, and constructs $q_0^{r(k-r)}$ distinct edges from graph planes over $\mathbf F_{q_0}$. \textsf{[JR2, JR3, JR8]} The remaining input is the construction of sufficiently many polynomial vectors and dual tests \emph{on each edge}, with their stated valuation and degree bounds. The source has separate treatments for $r\ge2$ and $r=1$; the latter also requires an averaged bound on $\kappa$. Neither quantitative construction is inferred merely from counting planes. \textsf{[JR9, JR10]}
